\section{Research gaps, recommended experiments and conclusions}
\subsection{What prior work already covers}
The search establishes that fluorescent conformal coatings on PCBs have already been inspected using UV candidate extraction, CNN verification, transfer learning, YOLO-style methods and Faster R-CNN \citep{lee2021cnn,lee2021transfer,choi2024uv,wang2025bubble,patel2026coverage}. More general PCB inspection has long used registration, reference difference and learned segmentation \citep{hara1988fluorescent,qub2025registration,ebayyeh2020aoi}. Commercial CCI systems already use coating plans, UV lighting and optional side cameras or thickness metrology \citep{goepelcci,nordsoncci,mirteccci}. These facts exclude simplistic novelty claims that AI, UV, reference masks, image tiling or hybrid pipelines are invented here.

\subsection{Research gaps defensible from the reviewed evidence}
\textbf{Gap 1: reproducible multi-defect evaluation.} The located direct studies generally target selected bubbles or coverage-related defect categories; comparable multi-class, board- and lot-disjoint error rates are not established by their public abstracts.\par
\textbf{Gap 2: observable scope and acceptance semantics.} The manufacturer's CAD-driven approach shows that a coating plan matters; an academic study can quantify the incremental value of explicit required/forbidden/occluded masks at matched false acceptance and rejection.\par
\textbf{Gap 3: illumination robustness under controlled cost.} UV intensity selection has already been studied \citep{choi2024uv}. What remains to test is whether reproducible correction and minimal additional illumination exposures reduce misses enough to offset line-time and hardware complexity.\par
\textbf{Gap 4: CPU edge validation.} Generic AI studies claim high GPU speed, but published results do not establish our complete cycle time on a cheap production-class Intel PC.\par
\textbf{Gap 5: source traceability and transferable datasets.} Existing open datasets cover bare PCB defects and general anomalies rather than a confirmed public, image-level UV coating benchmark \citep{lv2024dspcb,bergmann2019mvtec,visa2022dataset}. A properly documented new dataset is justified; publication rights are a separate requirement.\par

These are \emph{hypothesized research opportunities} relative to the currently reviewed sources, not proof that no other group has addressed them. Full-text searches of licensed proceedings and patents should precede any definitive novelty declaration.

\subsection{Recommended technical build sequence}
\begin{enumerate}[leftmargin=*]
\item \textbf{Optical feasibility bench.} Test the actual coating under 365~nm UV, verify the emission band with appropriate blocking, and quantify field uniformity using one camera. Acquire coated/uncoated references and carefully validated physical defects.
\item \textbf{Classical baseline first.} Implement image registration and flat-field correction, thresholding, morphological region metrics and the approved CAD coverage map. This provides immediate industrial value and an interpretable research benchmark.
\item \textbf{Representative dataset.} Expand to distinct physical boards and batches; label missing coating, overspray/keep-out violations, bubbles and dewetting only where experts can verify each category. Record the uninspectable surfaces.
\item \textbf{Lightweight AI comparison.} Train a small U-Net and one transformer alternative only after defining the baseline failures. Focus on rare, irregular or weakly fluorescent defects, not cosmetic improvement in already simple regions.
\item \textbf{Controlled hybrid comparison.} Test learned masking with and without reference constraints, correction and optional additional exposure. Report causal ablation results and miss cases.
\item \textbf{Edge and industrial qualification.} Export the chosen model to ONNX/OpenVINO, test multiple low-cost CPU systems, measure all timing stages and qualify the procedure against approved production acceptance thresholds.
\end{enumerate}

\subsection{Publication strategy}
A defensible scientific title would be \emph{``Reference-Aware UV Fluorescence Inspection of Conformal-Coated PCBs: A Controlled Comparison of Classical Vision and Lightweight Segmentation''}. Other viable emphases are low-data coating segmentation, multi-angle optical robustness or cost--sensitivity trade-offs. The most compelling paper should present a \emph{specific measured improvement} rather than an untested multi-component system description. Required figures include physical UV enclosure, emission/filter response, calibration map, defect taxonomy with expert labels, architecture, ablation performance curves, size-stratified misses and CPU cycle-time distributions. A release-approved evaluation subset and configuration-controlled source code would improve reproducibility; confidential production images need not be published to establish the methodology.

\subsection{Final synthesis}
For a 70~mm square assembly, controlled UV fluorescence and a single industrial camera offer a plausible economical path to optical verification of \emph{visible} coating coverage. The literature strongly supports building a reference-aware classical baseline before investing in larger AI models. It also demonstrates that the most dangerous research error would be to claim coating thickness, rare-defect reliability or production-ready speed based solely on a bright image or a generic benchmark. The contribution worth testing is an inspection \emph{system of measured optical signals, validated regions, independently verified defects, interpretable decisions and repeatable end-to-end timing}. Until those experiments exist, the present document should be used as a rigorous methodology and informed literature foundation, not as a results paper.
